\documentclass[11pt,a4paper]{article}

\usepackage[utf8]{inputenc}
\usepackage[T1]{fontenc}
\usepackage{lmodern}
\usepackage[margin=1in]{geometry}
\usepackage{hyperref}
\usepackage{booktabs}
\usepackage{longtable}
\usepackage{tabularx}
\usepackage{xcolor}
\usepackage{enumitem}
\usepackage{amsmath,amssymb}
\usepackage{graphicx}
\usepackage{titlesec}
\usepackage{fancyhdr}
\usepackage{array}

\definecolor{primaryblue}{RGB}{0, 51, 102}
\definecolor{secondaryblue}{RGB}{30, 90, 150}
\definecolor{darkgray}{RGB}{60, 60, 60}

\hypersetup{
    colorlinks=true,
    linkcolor=primaryblue,
    citecolor=primaryblue,
    urlcolor=secondaryblue,
    pdfauthor={ScholarMaster Research Initiative},
    pdftitle={ScholarMaster Master Research Portfolio Plan (P1-P25)}
}

\pagestyle{fancy}
\fancyhf{}
\fancyhead[L]{\small\textcolor{darkgray}{\textbf{ScholarMaster Research Initiative} --- Master Paper Plan}}
\fancyhead[R]{\small\textcolor{darkgray}{P1--P25 Governance Engine}}
\fancyfoot[C]{\thepage}
\renewcommand{\headrulewidth}{0.4pt}
\renewcommand{\footrulewidth}{0.4pt}

\titleformat{\section}{\Large\bfseries\color{primaryblue}}{\thesection}{1em}{}[\vspace{0.2em}\hrule\vspace{0.5em}]
\titleformat{\subsection}{\large\bfseries\color{secondaryblue}}{\thesubsection}{1em}{}
\titleformat{\subsubsection}{\normalsize\bfseries\color{darkgray}}{\thesubsubsection}{1em}{}

\newcolumntype{L}[1]{>{\raggedright\arraybackslash}p{#1}}
\newcolumntype{C}[1]{>{\centering\arraybackslash}p{#1}}
\newcolumntype{R}[1]{>{\raggedleft\arraybackslash}p{#1}}

\begin{document}

\begin{titlepage}
    \centering
    \vspace*{1.5cm}
    {\Huge \textbf{\color{primaryblue} SCHOLARMASTER}}\\[0.6cm]
    {\LARGE \textbf{Master Research Portfolio Plan}}\\[0.3cm]
    {\Large \textbf{Papers P1 through P25}}\\[1.5cm]
    \rule{\linewidth}{1.5pt}\\[0.6cm]
    {\Large \textbf{Research Architecture, Publication Strategy, Evidence Governance, Reviewer Calibration, and Future-Paper Framework}}\\[0.4cm]
    \rule{\linewidth}{1.5pt}\\[2.0cm]
    \begin{minipage}{0.85\textwidth}
        \centering
        \textbf{Authoritative Governance Engine Generated Strategic Roadmap}\\
        \vspace{0.4cm}
        \textit{Derived from Canonical Portfolio Registries and Calibrated Against Reviewer-6 Standard}\\
        \vspace{0.4cm}
        \textbf{Date}: August 29, 2026\\
        \textbf{Derived Governance Status}: CONDITIONAL (Companion dependency alerts or unverified items pending review)
    \end{minipage}
    \vfill
    {\small Swarnandhra College of Engineering \& Technology (Autonomous) $\cdot$ Edge-Native AI Systems Research Group}
\end{titlepage}

\newpage
\tableofcontents
\newpage

\section*{Executive Summary}
\addcontentsline{toc}{section}{Executive Summary}

The \textbf{ScholarMaster Research Initiative} represents a unified, multi-year scientific investigation into privacy-first, edge-native cyber-physical intelligence for academic and institutional environments. Across 25 distinct research manuscripts (\textbf{P1--P25}), the portfolio addresses the foundational tension between continuous, context-aware operational analytics and absolute individual privacy preservation.

This Master Paper Plan serves as the authoritative, source-driven strategic roadmap consolidating all planning, architectural formalisms, empirical validations, publication chronology audits, and reviewer calibration standards developed across the research program.

\section{Portfolio Overview and Structural Organization}
The portfolio spans four vertical strata governed by Constraint-First Architectural Synthesis (CFAS).

\section{Master Research Architecture and Invariant Governance}
Structural invariants compiled into shared memory enforce zero cross-layer memory or semantic leakage.

\section{Reviewer Calibration Standard (Paper-6 Framework)}
All manuscripts are calibrated against the 4 Reviewer-6 Skepticism Pillars (Novelty, Breadth, Language, Limitations).

\section{Publication and Citation Governance Policy}
Strict publication chronology is enforced: only verified published works (P5) and accepted/in-press works (P6) are cited.

% ==========================================
% INDIVIDUAL PAPER PLANS
% ==========================================

\section{Paper 1: ScholarMaster: A Layered Edge-Native Architecture for Real-Time Context-Aware Intelligent Systems}

\subsection*{Paper Identity \& Strategic Metadata}
\begin{itemize}[leftmargin=1.5em, itemsep=0.2em]
    \item \textbf{Paper Identifier}: \texttt{P01} $\cdot$ \textbf{Methodological Type}: SYSTEMS
    \item \textbf{Research Area}: Edge Computing / Systems Architecture
    \item \textbf{Publication Status}: \textbf{DRAFT}
    \item \textbf{Target Venues}: Primary: \textit{IEEE Transactions on Edge Computing} $\cdot$ Alternatives: \textit{ACM Trans. Cyber-Physical Systems, IEEE Internet of Things Journal}
\end{itemize}

\subsection*{Scientific & Architectural Contribution Decomposition}
\begin{itemize}[leftmargin=1.5em, itemsep=0.2em]
    \item \textbf{Known Primitives Acknowledged}: POSIX shared memory ring buffers, multi-threaded IPC, microservices.
    \item \textbf{Unique Subsystem Contribution}: 4-Stratum architectural reference model with non-bypassable boundary invariants preventing cross-layer semantic and memory leakage.
    \item \textbf{Novelty Claim Formulation}: Formalized layered decomposition and volatile write-once boundary enforcement.
    \item \textbf{Theoretical Formalisms}: No formal mathematical theorem blocks.
\end{itemize}

\subsection*{Empirical Evidence & Telemetry}
\begin{itemize}[leftmargin=1.5em, itemsep=0.2em]
    \item \textbf{Key Numerical Telemetry}: 0.42 ms to 1.1 ms inter-stratum transfer latency (Source: Table IV empirical latency measurements across 10,000 trace events)
    \item \textbf{Remaining Reviewer-6 Risk}: Engineering integration vs new computational primitive (mitigated via explicit scoping to cyber-physical reference stack).
\end{itemize}

\section{Paper 2: A Context-Aware Multi-Modal Framework for Asymmetric Risk Control in Student Engagement Analysis}

\subsection*{Paper Identity \& Strategic Metadata}
\begin{itemize}[leftmargin=1.5em, itemsep=0.2em]
    \item \textbf{Paper Identifier}: \texttt{P02} $\cdot$ \textbf{Methodological Type}: HYBRID (EMPIRICAL / MULTIMODAL)
    \item \textbf{Research Area}: Affective Computing / Multimodal AI
    \item \textbf{Publication Status}: \textbf{DRAFT}
    \item \textbf{Target Venues}: Primary: \textit{IEEE Transactions on Learning Technologies} $\cdot$ Alternatives: \textit{ACM Trans. Interactive Intelligent Systems, Computers \& Education}
\end{itemize}

\subsection*{Scientific & Architectural Contribution Decomposition}
\begin{itemize}[leftmargin=1.5em, itemsep=0.2em]
    \item \textbf{Known Primitives Acknowledged}: Bayesian cost-sensitive risk minimization, IIR temporal smoothing, facial expression valence models.
    \item \textbf{Unique Subsystem Contribution}: Context-aware asymmetric loss optimization dynamically re-weighting multi-modal streams during high-load STEM tasks.
    \item \textbf{Novelty Claim Formulation}: Theorem 1 asymmetric risk bound minimizing Type-II false negative engagement classifications.
    \item \textbf{Theoretical Formalisms}: 2 formal mathematical blocks (proposition, theorem)
\end{itemize}

\subsection*{Empirical Evidence & Telemetry}
\begin{itemize}[leftmargin=1.5em, itemsep=0.2em]
    \item \textbf{Key Numerical Telemetry}: Reduction in Type-II false negatives from 42.0\% to 6.0\% (Source: Sim-Class-24 simulation benchmark harness (N=1,440 session traces))
    \item \textbf{Remaining Reviewer-6 Risk}: Synthetic simulation harness (Sim-Class-24) vs unconstrained physical classroom acoustics.
\end{itemize}

\section{Paper 3: Pose-Only Edge Action Sensing with Enforced Volatile Memory Confinement}

\subsection*{Paper Identity \& Strategic Metadata}
\begin{itemize}[leftmargin=1.5em, itemsep=0.2em]
    \item \textbf{Paper Identifier}: \texttt{P03} $\cdot$ \textbf{Methodological Type}: HARDWARE/EDGE (VISION)
    \item \textbf{Research Area}: Edge Computer Vision / Privacy Engineering
    \item \textbf{Publication Status}: \textbf{DRAFT}
    \item \textbf{Target Venues}: Primary: \textit{IEEE Transactions on Circuits and Systems for Video Technology} $\cdot$ Alternatives: \textit{IEEE Trans. Information Forensics and Security, ACM Trans. Multimedia}
\end{itemize}

\subsection*{Scientific & Architectural Contribution Decomposition}
\begin{itemize}[leftmargin=1.5em, itemsep=0.2em]
    \item \textbf{Known Primitives Acknowledged}: 17-keypoint human pose estimation (MediaPipe/OpenPose), lightweight GCN action recognition.
    \item \textbf{Unique Subsystem Contribution}: Volatile-only memory ring buffer zeroization with metric-space dimensional collapse proof.
    \item \textbf{Novelty Claim Formulation}: Theorem 1 exact metric pixel non-invertibility under enforced memory destruction.
    \item \textbf{Theoretical Formalisms}: 1 formal mathematical blocks (theorem)
\end{itemize}

\subsection*{Empirical Evidence & Telemetry}
\begin{itemize}[leftmargin=1.5em, itemsep=0.2em]
    \item \textbf{Key Numerical Telemetry}: Documented in manuscript.
    \item \textbf{Remaining Reviewer-6 Risk}: Generative diffusion pose-to-image hallucination (addressed by explicitly bounding threat model).
\end{itemize}

\section{Paper 4: Real-Time Schedule Compliance via Spatiotemporal Predicate Evaluation and Relational Lookup}

\subsection*{Paper Identity \& Strategic Metadata}
\begin{itemize}[leftmargin=1.5em, itemsep=0.2em]
    \item \textbf{Paper Identifier}: \texttt{P04} $\cdot$ \textbf{Methodological Type}: SYSTEMS (STREAM / DATABASE)
    \item \textbf{Research Area}: Cyber-Physical Systems / Database Stream Processing
    \item \textbf{Publication Status}: \textbf{DRAFT}
    \item \textbf{Target Venues}: Primary: \textit{ACM Transactions on Cyber-Physical Systems} $\cdot$ Alternatives: \textit{IEEE Trans. Services Computing, IEEE Internet of Things Journal}
\end{itemize}

\subsection*{Scientific & Architectural Contribution Decomposition}
\begin{itemize}[leftmargin=1.5em, itemsep=0.2em]
    \item \textbf{Known Primitives Acknowledged}: Relational connection pooling, list-based PostgreSQL table partitioning, temporal debounce filters.
    \item \textbf{Unique Subsystem Contribution}: Continuous Predicate Evaluation Model (CPEM) with Probabilistic Cumulative Violation Filter (PCVF) for burst-elastic compliance.
    \item \textbf{Novelty Claim Formulation}: Theorem 1 transient violation suppression invariant and Theorem 2 bounded relational lookup latency.
    \item \textbf{Theoretical Formalisms}: 2 formal mathematical blocks (theorem)
\end{itemize}

\subsection*{Empirical Evidence & Telemetry}
\begin{itemize}[leftmargin=1.5em, itemsep=0.2em]
    \item \textbf{Key Numerical Telemetry}: 5,000 QPS with p99 latency < 15 ms (Source: PostgreSQL connection pooling load test (50,000 synthetic events))
    \item \textbf{Remaining Reviewer-6 Risk}: Debounce filter is a discrete leaky-bucket (acknowledged in Section I-D).
\end{itemize}

\section{Paper 5: Memory-Bound Edge Efficiency Envelope (MBEEE): A Hardware-Level Analytical Model}

\subsection*{Paper Identity \& Strategic Metadata}
\begin{itemize}[leftmargin=1.5em, itemsep=0.2em]
    \item \textbf{Paper Identifier}: \texttt{P05} $\cdot$ \textbf{Methodological Type}: HARDWARE/EDGE (ANALYTICAL MODEL)
    \item \textbf{Research Area}: Embedded Systems / Hardware Modeling
    \item \textbf{Publication Status}: \textbf{PUBLISHED}
    \item \textbf{Target Venues}: Primary: \textit{IEEE Access (PUBLISHED BASELINE)} $\cdot$ Alternatives: \textit{N/A (Published), N/A}
\end{itemize}

\subsection*{Scientific & Architectural Contribution Decomposition}
\begin{itemize}[leftmargin=1.5em, itemsep=0.2em]
    \item \textbf{Known Primitives Acknowledged}: Unified Memory Architecture (UMA), PCIe interconnects, Roofline model.
    \item \textbf{Unique Subsystem Contribution}: Hardware-level analytical model proving zero PCIe bus overhead for edge vision pipelines.
    \item \textbf{Novelty Claim Formulation}: PUBLISHED BASELINE (IEEE Access 2026).
    \item \textbf{Theoretical Formalisms}: No formal mathematical theorem blocks.
\end{itemize}

\subsection*{Empirical Evidence & Telemetry}
\begin{itemize}[leftmargin=1.5em, itemsep=0.2em]
    \item \textbf{Key Numerical Telemetry}: Documented in manuscript.
    \item \textbf{Remaining Reviewer-6 Risk}: None (Published manuscript).
\end{itemize}

\section{Paper 6: NLOS Acoustic Sensing via Spectral Gating and GCC-PHAT}

\subsection*{Paper Identity \& Strategic Metadata}
\begin{itemize}[leftmargin=1.5em, itemsep=0.2em]
    \item \textbf{Paper Identifier}: \texttt{P06} $\cdot$ \textbf{Methodological Type}: HARDWARE/EDGE (ACOUSTIC SIGNAL PROCESSING)
    \item \textbf{Research Area}: Acoustic Signal Processing / Anomaly Detection
    \item \textbf{Publication Status}: \textbf{ACCEPTED}
    \item \textbf{Target Venues}: Primary: \textit{ACCEPTED / IN PRESS BASELINE} $\cdot$ Alternatives: \textit{N/A (Accepted), N/A}
\end{itemize}

\subsection*{Scientific & Architectural Contribution Decomposition}
\begin{itemize}[leftmargin=1.5em, itemsep=0.2em]
    \item \textbf{Known Primitives Acknowledged}: Spectral gating, GCC-PHAT time-difference-of-arrival, acoustic impulse responses.
    \item \textbf{Unique Subsystem Contribution}: Non-line-of-sight acoustic anomaly detection with privacy-preserving feature extraction.
    \item \textbf{Novelty Claim Formulation}: ACCEPTED BASELINE (Accepted / In Press).
    \item \textbf{Theoretical Formalisms}: No formal mathematical theorem blocks.
\end{itemize}

\subsection*{Empirical Evidence & Telemetry}
\begin{itemize}[leftmargin=1.5em, itemsep=0.2em]
    \item \textbf{Key Numerical Telemetry}: Documented in manuscript.
    \item \textbf{Remaining Reviewer-6 Risk}: Single-corridor validation (addressed during peer review).
\end{itemize}

\section{Paper 7: Sub-Millisecond Identity Retrieval via HNSW + LDCC}

\subsection*{Paper Identity \& Strategic Metadata}
\begin{itemize}[leftmargin=1.5em, itemsep=0.2em]
    \item \textbf{Paper Identifier}: \texttt{P07} $\cdot$ \textbf{Methodological Type}: SYSTEMS (RETRIEVAL ALGORITHM)
    \item \textbf{Research Area}: Information Retrieval / Biometric Vector Search
    \item \textbf{Publication Status}: \textbf{DRAFT}
    \item \textbf{Target Venues}: Primary: \textit{IEEE Transactions on Information Forensics and Security} $\cdot$ Alternatives: \textit{IEEE Trans. Pattern Analysis and Machine Intelligence, ACM Trans. Information Systems}
\end{itemize}

\subsection*{Scientific & Architectural Contribution Decomposition}
\begin{itemize}[leftmargin=1.5em, itemsep=0.2em]
    \item \textbf{Known Primitives Acknowledged}: Hierarchical Navigable Small World (HNSW) graph indexing, k-NN local density estimation.
    \item \textbf{Unique Subsystem Contribution}: Integration of Local Density Consistency Criterion (LDCC) into HNSW graph frontier traversals for sub-millisecond open-set unknown rejection.
    \item \textbf{Novelty Claim Formulation}: Theorem 1 logarithmic traversal scaling and Theorem 2 LDCC open-set unknown rejection bound on pre-extracted embeddings.
    \item \textbf{Theoretical Formalisms}: 2 formal mathematical blocks (theorem)
\end{itemize}

\subsection*{Empirical Evidence & Telemetry}
\begin{itemize}[leftmargin=1.5em, itemsep=0.2em]
    \item \textbf{Key Numerical Telemetry}: 0.72 ms query latency on N=100,000 vector gallery (Source: HNSW (M=16, efSearch=64) + LDCC retrieval on embedded CPU)
    \item \textbf{Remaining Reviewer-6 Risk}: Requires pre-extracted embeddings (explicitly scoped in Abstract and Section I).
\end{itemize}

\section{Paper 8: A Cryptographic Provenance Model with Erasure-Compatible Immutability}

\subsection*{Paper Identity \& Strategic Metadata}
\begin{itemize}[leftmargin=1.5em, itemsep=0.2em]
    \item \textbf{Paper Identifier}: \texttt{P08} $\cdot$ \textbf{Methodological Type}: SECURITY (APPLIED CRYPTOGRAPHY)
    \item \textbf{Research Area}: Applied Cryptography / Privacy Engineering
    \item \textbf{Publication Status}: \textbf{DRAFT}
    \item \textbf{Target Venues}: Primary: \textit{IEEE Transactions on Dependable and Secure Computing} $\cdot$ Alternatives: \textit{ACM Trans. Privacy and Security, IEEE Trans. Information Forensics and Security}
\end{itemize}

\subsection*{Scientific & Architectural Contribution Decomposition}
\begin{itemize}[leftmargin=1.5em, itemsep=0.2em]
    \item \textbf{Known Primitives Acknowledged}: Merkle trees, cryptographic key shredding, permissioned blockchain ledgers.
    \item \textbf{Unique Subsystem Contribution}: Per-Identity Symmetric Key (PISK) architecture enabling GDPR Article 17 erasure without invalidating historical Merkle tree roots.
    \item \textbf{Novelty Claim Formulation}: Erasure-compatible immutable provenance architecture for high-frequency edge telemetry.
    \item \textbf{Theoretical Formalisms}: No formal mathematical theorem blocks.
\end{itemize}

\subsection*{Empirical Evidence & Telemetry}
\begin{itemize}[leftmargin=1.5em, itemsep=0.2em]
    \item \textbf{Key Numerical Telemetry}: 1,200 TPS throughput on permissioned ledger (Source: Batched Merkle tree benchmark across N=100,000 transaction events)
    \item \textbf{Remaining Reviewer-6 Risk}: KMS root custody trust dependency (explicitly stated in Section I-D).
\end{itemize}

\section{Paper 9: A Hierarchical Edge Control Plane for Policy-Aware Multi-Module AI Orchestration}

\subsection*{Paper Identity \& Strategic Metadata}
\begin{itemize}[leftmargin=1.5em, itemsep=0.2em]
    \item \textbf{Paper Identifier}: \texttt{P09} $\cdot$ \textbf{Methodological Type}: SYSTEMS (CONTROL PLANE)
    \item \textbf{Research Area}: Edge Computing / Real-Time Control
    \item \textbf{Publication Status}: \textbf{DRAFT}
    \item \textbf{Target Venues}: Primary: \textit{IEEE Transactions on Mobile Computing} $\cdot$ Alternatives: \textit{ACM Trans. Embedded Computing Systems, IEEE Internet of Things Journal}
\end{itemize}

\subsection*{Scientific & Architectural Contribution Decomposition}
\begin{itemize}[leftmargin=1.5em, itemsep=0.2em]
    \item \textbf{Known Primitives Acknowledged}: Lyapunov queue stability, dynamic duty cycling, state machine schedulers.
    \item \textbf{Unique Subsystem Contribution}: Kinematic-coupled sampling bound and Lyapunov-stable rate governor for heterogeneous edge AI perception.
    \item \textbf{Novelty Claim Formulation}: Theorem 1 kinematic sampling bound and Theorem 2 Lyapunov rate stability.
    \item \textbf{Theoretical Formalisms}: 2 formal mathematical blocks (theorem)
\end{itemize}

\subsection*{Empirical Evidence & Telemetry}
\begin{itemize}[leftmargin=1.5em, itemsep=0.2em]
    \item \textbf{Key Numerical Telemetry}: Documented in manuscript.
    \item \textbf{Remaining Reviewer-6 Risk}: Phase transitions assume predictable schedule blocks (fail-safe fallback detailed in Section X).
\end{itemize}

\section{Paper 10: Integrated Stress Validation of an Edge-Native Academic Analytics Architecture}

\subsection*{Paper Identity \& Strategic Metadata}
\begin{itemize}[leftmargin=1.5em, itemsep=0.2em]
    \item \textbf{Paper Identifier}: \texttt{P10} $\cdot$ \textbf{Methodological Type}: EMPIRICAL (LONGITUDINAL STRESS BENCHMARK)
    \item \textbf{Research Area}: Software Reliability / Empirical Systems Engineering
    \item \textbf{Publication Status}: \textbf{DRAFT}
    \item \textbf{Target Venues}: Primary: \textit{IEEE Transactions on Reliability} $\cdot$ Alternatives: \textit{Empirical Software Engineering, IEEE Systems Journal}
\end{itemize}

\subsection*{Scientific & Architectural Contribution Decomposition}
\begin{itemize}[leftmargin=1.5em, itemsep=0.2em]
    \item \textbf{Known Primitives Acknowledged}: Software reliability testing, fault injection harnesses, burn-in protocols.
    \item \textbf{Unique Subsystem Contribution}: Comprehensive 168-hour integrated stress matrix evaluation on physical ARM64 hardware under compound stressors.
    \item \textbf{Novelty Claim Formulation}: Empirical systems resilience characterization across 7 days of continuous unmonitored execution.
    \item \textbf{Theoretical Formalisms}: No formal mathematical theorem blocks.
\end{itemize}

\subsection*{Empirical Evidence & Telemetry}
\begin{itemize}[leftmargin=1.5em, itemsep=0.2em]
    \item \textbf{Key Numerical Telemetry}: 168 hours continuous autonomous execution without memory leak (Source: Physical 7-day burn-in telemetry log on ARM64 Linux appliance)
    \item \textbf{Remaining Reviewer-6 Risk}: Empirical evaluation without new algorithmic theorems (framed honestly as systems research).
\end{itemize}

\section{Paper 11: Lifecycle Hardening of Immutable Edge Appliances under Power and Connectivity Instability}

\subsection*{Paper Identity \& Strategic Metadata}
\begin{itemize}[leftmargin=1.5em, itemsep=0.2em]
    \item \textbf{Paper Identifier}: \texttt{P11} $\cdot$ \textbf{Methodological Type}: SYSTEMS (EMBEDDED OS LIFECYCLE)
    \item \textbf{Research Area}: Embedded Operating Systems / System Hardening
    \item \textbf{Publication Status}: \textbf{DRAFT}
    \item \textbf{Target Venues}: Primary: \textit{IEEE Transactions on Industrial Informatics} $\cdot$ Alternatives: \textit{ACM Trans. Embedded Computing Systems, IEEE Trans. Dependable and Secure Computing}
\end{itemize}

\subsection*{Scientific & Architectural Contribution Decomposition}
\begin{itemize}[leftmargin=1.5em, itemsep=0.2em]
    \item \textbf{Known Primitives Acknowledged}: Squashfs, overlayfs, dual A/B partitioning, U-Boot watchdog rollback.
    \item \textbf{Unique Subsystem Contribution}: State-invariance proof (Theorem 1) and bounded rollback liveness (Lemma 1) for unattended edge appliances under unannounced power loss.
    \item \textbf{Novelty Claim Formulation}: Formalized immutable lifecycle stack and 50-cycle physical power-cut validation.
    \item \textbf{Theoretical Formalisms}: 2 formal mathematical blocks (theorem, lemma)
\end{itemize}

\subsection*{Empirical Evidence & Telemetry}
\begin{itemize}[leftmargin=1.5em, itemsep=0.2em]
    \item \textbf{Key Numerical Telemetry}: 50 physical power-cut cycles with zero state corruption (Source: Physical hardware power-cut test rig with A/B squashfs/overlayfs rollback)
    \item \textbf{Remaining Reviewer-6 Risk}: Standard embedded Linux components unified under formal invariance proofs.
\end{itemize}

\section{Paper 12: Flash Endurance Engineering for Write-Intensive Embedded Workloads: Extending SD Card Lifespan Through Kernel-Level Optimization}

\subsection*{Paper Identity \& Strategic Metadata}
\begin{itemize}[leftmargin=1.5em, itemsep=0.2em]
    \item \textbf{Paper Identifier}: \texttt{P12} $\cdot$ \textbf{Methodological Type}: SYSTEMS (STORAGE OPTIMIZATION)
    \item \textbf{Research Area}: Storage Systems / Embedded Linux
    \item \textbf{Publication Status}: \textbf{DRAFT}
    \item \textbf{Target Venues}: Primary: \textit{ACM Transactions on Storage} $\cdot$ Alternatives: \textit{IEEE Trans. Computer-Aided Design of Integrated Circuits, IEEE Embedded Systems Letters}
\end{itemize}

\subsection*{Scientific & Architectural Contribution Decomposition}
\begin{itemize}[leftmargin=1.5em, itemsep=0.2em]
    \item \textbf{Known Primitives Acknowledged}: F2FS, Linux tmpfs, VFS dirty-ratio sysctl tuning, Flash Translation Layer (FTL) wear leveling.
    \item \textbf{Unique Subsystem Contribution}: End-to-end kernel-to-VFS optimization stack reducing Write Amplification Factor (WAF) from 3.42 to 1.12 on flash memory.
    \item \textbf{Novelty Claim Formulation}: Systematic characterization of edge AI write amplification and analytical 30x lifespan extension model.
    \item \textbf{Theoretical Formalisms}: No formal mathematical theorem blocks.
\end{itemize}

\subsection*{Empirical Evidence & Telemetry}
\begin{itemize}[leftmargin=1.5em, itemsep=0.2em]
    \item \textbf{Key Numerical Telemetry}: 30x analytical lifespan extension projection / WAF reduction from 3.42 to 1.12 (Source: Analytical NAND wear-out model evaluated on 24-hour blktrace empirical write volumes)
    \item \textbf{Remaining Reviewer-6 Risk}: Lifespan extension is an analytical projection from WAF models (explicitly qualified in text).
\end{itemize}

\section{Paper 13: Federated Drift Compensation via Active Learning in Edge Deployed Neural Networks}

\subsection*{Paper Identity \& Strategic Metadata}
\begin{itemize}[leftmargin=1.5em, itemsep=0.2em]
    \item \textbf{Paper Identifier}: \texttt{P13} $\cdot$ \textbf{Methodological Type}: DISTRIBUTED ML (ACTIVE FEDERATED LEARNING)
    \item \textbf{Research Area}: Distributed Machine Learning / Privacy-Preserving AI
    \item \textbf{Publication Status}: \textbf{DRAFT}
    \item \textbf{Target Venues}: Primary: \textit{IEEE Transactions on Neural Networks and Learning Systems} $\cdot$ Alternatives: \textit{IEEE Trans. Information Forensics and Security, ACM Trans. Intelligent Systems and Technology}
\end{itemize}

\subsection*{Scientific & Architectural Contribution Decomposition}
\begin{itemize}[leftmargin=1.5em, itemsep=0.2em]
    \item \textbf{Known Primitives Acknowledged}: Bayesian Active Learning by Disagreement (BALD), Differential Privacy (DP), Federated Averaging.
    \item \textbf{Unique Subsystem Contribution}: Stationary variance bound proof (Theorem 1) for federated active learning under differentially private gradient perturbation.
    \item \textbf{Novelty Claim Formulation}: Theorem 1 stationary variance bound and 93.0\% convergence under 15\% annotation budget.
    \item \textbf{Theoretical Formalisms}: 1 formal mathematical blocks (theorem)
\end{itemize}

\subsection*{Empirical Evidence & Telemetry}
\begin{itemize}[leftmargin=1.5em, itemsep=0.2em]
    \item \textbf{Key Numerical Telemetry}: 93.0\% accuracy under 15\% annotation budget (Source: Simulated 10-node federated active learning cluster with BALD acquisition)
    \item \textbf{Remaining Reviewer-6 Risk}: Simulated federated node topology (explicitly qualified in Abstract and Section VI).
\end{itemize}

\section{Paper 14: Hierarchical Federated Aggregation for Cross-Institution Model Adaptation Under Asynchronous Participation}

\subsection*{Paper Identity \& Strategic Metadata}
\begin{itemize}[leftmargin=1.5em, itemsep=0.2em]
    \item \textbf{Paper Identifier}: \texttt{P14} $\cdot$ \textbf{Methodological Type}: DISTRIBUTED ML (HIERARCHICAL OPTIMIZATION)
    \item \textbf{Research Area}: Distributed Optimization / Federated Learning
    \item \textbf{Publication Status}: \textbf{DRAFT}
    \item \textbf{Target Venues}: Primary: \textit{IEEE Transactions on Parallel and Distributed Systems} $\cdot$ Alternatives: \textit{IEEE Trans. Big Data, ACM Trans. Autonomous and Adaptive Systems}
\end{itemize}

\subsection*{Scientific & Architectural Contribution Decomposition}
\begin{itemize}[leftmargin=1.5em, itemsep=0.2em]
    \item \textbf{Known Primitives Acknowledged}: Hierarchical federated learning (HFL), asynchronous FedAvg, Dirichlet label skew.
    \item \textbf{Unique Subsystem Contribution}: Topology-aware polynomial staleness dampening with asymptotic convergence proof under bounded staleness.
    \item \textbf{Novelty Claim Formulation}: Theorem 1 asymptotic convergence under dual-level non-IID divergence and asynchronous stragglers.
    \item \textbf{Theoretical Formalisms}: 1 formal mathematical blocks (theorem)
\end{itemize}

\subsection*{Empirical Evidence & Telemetry}
\begin{itemize}[leftmargin=1.5em, itemsep=0.2em]
    \item \textbf{Key Numerical Telemetry}: Documented in manuscript.
    \item \textbf{Remaining Reviewer-6 Risk}: Simulated multi-institutional WAN network (acknowledged in Section VIII).
\end{itemize}

\section{Paper 15: Augmented Situation Awareness: Reducing Cognitive Load in Campus Security via Spatially-Anchored AR Visualization}

\subsection*{Paper Identity \& Strategic Metadata}
\begin{itemize}[leftmargin=1.5em, itemsep=0.2em]
    \item \textbf{Paper Identifier}: \texttt{P15} $\cdot$ \textbf{Methodological Type}: HCI/USER STUDY (SPATIAL AR)
    \item \textbf{Research Area}: Human-Computer Interaction / Spatial Computing
    \item \textbf{Publication Status}: \textbf{DRAFT}
    \item \textbf{Target Venues}: Primary: \textit{IEEE Transactions on Visualization and Computer Graphics} $\cdot$ Alternatives: \textit{ACM CHI Conference on Human Factors in Computing Systems, IEEE ISMAR}
\end{itemize}

\subsection*{Scientific & Architectural Contribution Decomposition}
\begin{itemize}[leftmargin=1.5em, itemsep=0.2em]
    \item \textbf{Known Primitives Acknowledged}: Visual Positioning System (VPS), SLAM spatial anchors, NASA-TLX cognitive load index.
    \item \textbf{Unique Subsystem Contribution}: Deterministic AR projection latency bound (Theorem 1) and empirical incident response user study.
    \item \textbf{Novelty Claim Formulation}: Theorem 1 projection latency bound and 34.2\% time-to-action reduction across 20 participants.
    \item \textbf{Theoretical Formalisms}: 2 formal mathematical blocks (proposition, theorem)
\end{itemize}

\subsection*{Empirical Evidence & Telemetry}
\begin{itemize}[leftmargin=1.5em, itemsep=0.2em]
    \item \textbf{Key Numerical Telemetry}: N=20 participants in controlled staged user study / 34.2\% TTA reduction (Source: 20-participant IRB-approved within-subject AR evaluation protocol)
    \item \textbf{Remaining Reviewer-6 Risk}: Controlled staged drills vs live emergency crisis panic (analyzed in Section IX).
\end{itemize}

\section{Paper 16: Beyond the Panopticon: A Longitudinal Study of Student Trust in Automated Stewardship Systems}

\subsection*{Paper Identity \& Strategic Metadata}
\begin{itemize}[leftmargin=1.5em, itemsep=0.2em]
    \item \textbf{Paper Identifier}: \texttt{P16} $\cdot$ \textbf{Methodological Type}: HCI/USER STUDY (LONGITUDINAL STS)
    \item \textbf{Research Area}: Human-Computer Interaction / Science \& Technology Studies (STS)
    \item \textbf{Publication Status}: \textbf{DRAFT}
    \item \textbf{Target Venues}: Primary: \textit{ACM CHI Conference on Human Factors in Computing Systems} $\cdot$ Alternatives: \textit{ACM CSCW, IEEE Technology and Society Magazine}
\end{itemize}

\subsection*{Scientific & Architectural Contribution Decomposition}
\begin{itemize}[leftmargin=1.5em, itemsep=0.2em]
    \item \textbf{Known Primitives Acknowledged}: System Trust Scale (STS), Panopticon theory (Foucault/Bentham), Mayer's trust model.
    \item \textbf{Unique Subsystem Contribution}: 16-week longitudinal quantification of student trust mechanics under visible architectural privacy cues (Skeleton View).
    \item \textbf{Novelty Claim Formulation}: Empirical quantification demonstrating that visible abstraction dominates backend policy claims in establishing human trust.
    \item \textbf{Theoretical Formalisms}: No formal mathematical theorem blocks.
\end{itemize}

\subsection*{Empirical Evidence & Telemetry}
\begin{itemize}[leftmargin=1.5em, itemsep=0.2em]
    \item \textbf{Key Numerical Telemetry}: N=312 students in 16-week longitudinal STS study / 68\% trust increase (Source: 16-week pre/post empirical survey across 4 academic departments)
    \item \textbf{Remaining Reviewer-6 Risk}: Single-institution cultural cohort (acknowledged in Section VI).
\end{itemize}

\section{Paper 17: Architectural Irreversibility: Enforcing Privacy, Governance, and Trust in Intelligent Campus Systems}

\subsection*{Paper Identity \& Strategic Metadata}
\begin{itemize}[leftmargin=1.5em, itemsep=0.2em]
    \item \textbf{Paper Identifier}: \texttt{P17} $\cdot$ \textbf{Methodological Type}: CONCEPTUAL (INVITED VISION DOCTRINE)
    \item \textbf{Research Area}: Position Paper / Privacy Architecture Doctrine
    \item \textbf{Publication Status}: \textbf{DRAFT}
    \item \textbf{Target Venues}: Primary: \textit{IEEE Security \& Privacy Magazine (INVITED POSITION TRACK)} $\cdot$ Alternatives: \textit{IEEE Computer, Communications of the ACM (Viewpoints)}
\end{itemize}

\subsection*{Scientific & Architectural Contribution Decomposition}
\begin{itemize}[leftmargin=1.5em, itemsep=0.2em]
    \item \textbf{Known Primitives Acknowledged}: Privacy-by-Design (Cavoukian), differential privacy concepts, institutional compliance.
    \item \textbf{Unique Subsystem Contribution}: Conceptual architectural doctrine formalizing Privacy-by-Architecture over Privacy-by-Policy.
    \item \textbf{Novelty Claim Formulation}: Foundational architectural position treatise on irreversible physical sensing.
    \item \textbf{Theoretical Formalisms}: No formal mathematical theorem blocks.
\end{itemize}

\subsection*{Empirical Evidence & Telemetry}
\begin{itemize}[leftmargin=1.5em, itemsep=0.2em]
    \item \textbf{Key Numerical Telemetry}: Documented in manuscript.
    \item \textbf{Remaining Reviewer-6 Risk}: Treatise format lacking empirical tables/proofs (explicitly positioned as Invited Vision Paper).
\end{itemize}

\section{Paper 18: Runtime Enforcement of Architectural Irreversibility in Edge AI Systems This research was conducted at the Department of Computer Science and Engineering, Swarnandhra College of Engineering \& Technology (Autonomous), as part of academic research activities.}

\subsection*{Paper Identity \& Strategic Metadata}
\begin{itemize}[leftmargin=1.5em, itemsep=0.2em]
    \item \textbf{Paper Identifier}: \texttt{P18} $\cdot$ \textbf{Methodological Type}: SYSTEMS (RUNTIME SECURITY VERIFICATION)
    \item \textbf{Research Area}: Systems Security / Runtime Verification
    \item \textbf{Publication Status}: \textbf{DRAFT}
    \item \textbf{Target Venues}: Primary: \textit{IEEE Transactions on Dependable and Secure Computing} $\cdot$ Alternatives: \textit{ACM Trans. Computer Systems, IEEE Trans. Information Forensics and Security}
\end{itemize}

\subsection*{Scientific & Architectural Contribution Decomposition}
\begin{itemize}[leftmargin=1.5em, itemsep=0.2em]
    \item \textbf{Known Primitives Acknowledged}: POSIX timerfd, SIGKILL signals, shared memory IPC, watchdog daemons.
    \item \textbf{Unique Subsystem Contribution}: Process-isolated runtime enforcement harness validating 100\% fail-closed state reachability across 6 fault scenarios with zero memory residue.
    \item \textbf{Novelty Claim Formulation}: Comprehensive runtime verification methodology and fail-closed state machine for edge AI nodes.
    \item \textbf{Theoretical Formalisms}: No formal mathematical theorem blocks.
\end{itemize}

\subsection*{Empirical Evidence & Telemetry}
\begin{itemize}[leftmargin=1.5em, itemsep=0.2em]
    \item \textbf{Key Numerical Telemetry}: Documented in manuscript.
    \item \textbf{Remaining Reviewer-6 Risk}: Relies on standard POSIX OS signals (framed as an applied runtime verification architecture).
\end{itemize}

\section{Paper 19: Formal Threat Model and Trusted Computing Base Definition for Architecturally Irreversible Edge AI This research was conducted at the Department of Computer Science and Engineering, Swarnandhra College of Engineering \& Technology (Autonomous), as part of academic research activities.}

\subsection*{Paper Identity \& Strategic Metadata}
\begin{itemize}[leftmargin=1.5em, itemsep=0.2em]
    \item \textbf{Paper Identifier}: \texttt{P19} $\cdot$ \textbf{Methodological Type}: THEORETICAL (FORMAL SECURITY TCB)
    \item \textbf{Research Area}: Formal Methods / Hardware Security
    \item \textbf{Publication Status}: \textbf{DRAFT}
    \item \textbf{Target Venues}: Primary: \textit{IEEE Transactions on Information Forensics and Security} $\cdot$ Alternatives: \textit{ACM Trans. Privacy and Security, IEEE Trans. Dependable and Secure Computing}
\end{itemize}

\subsection*{Scientific & Architectural Contribution Decomposition}
\begin{itemize}[leftmargin=1.5em, itemsep=0.2em]
    \item \textbf{Known Primitives Acknowledged}: Saltzer-Schroeder protection, Goguen-Meseguer non-interference, Metric Temporal Logic (MTL).
    \item \textbf{Unique Subsystem Contribution}: Minimal TCB definition and formal proof of 5 foundational security theorems for irreversible edge AI.
    \item \textbf{Novelty Claim Formulation}: Theorems 1-5 (Pixel Non-Reconstructability, Non-Interference, Temporal Exposure, Governance Gate, Safe Halt).
    \item \textbf{Theoretical Formalisms}: 5 formal mathematical blocks (theorem)
\end{itemize}

\subsection*{Empirical Evidence & Telemetry}
\begin{itemize}[leftmargin=1.5em, itemsep=0.2em]
    \item \textbf{Key Numerical Telemetry}: Documented in manuscript.
    \item \textbf{Remaining Reviewer-6 Risk}: Theoretical security model vs physical side-channel probing (explicitly bounded in Section VII).
\end{itemize}

\section{Paper 20: The ScholarMaster Architecture: A Unified Reference Model for Privacy-First Intelligent Campus Systems}

\subsection*{Paper Identity \& Strategic Metadata}
\begin{itemize}[leftmargin=1.5em, itemsep=0.2em]
    \item \textbf{Paper Identifier}: \texttt{P20} $\cdot$ \textbf{Methodological Type}: ARCHITECTURAL (META-SYSTEMS SURVEY)
    \item \textbf{Research Area}: Reference Architecture / Meta-Systems Survey
    \item \textbf{Publication Status}: \textbf{DRAFT}
    \item \textbf{Target Venues}: Primary: \textit{IEEE Communications Surveys \& Tutorials} $\cdot$ Alternatives: \textit{ACM Computing Surveys, IEEE Systems Journal}
\end{itemize}

\subsection*{Scientific & Architectural Contribution Decomposition}
\begin{itemize}[leftmargin=1.5em, itemsep=0.2em]
    \item \textbf{Known Primitives Acknowledged}: EdgeX Foundry, NIST CPS framework, microservices, 4-stratum computing.
    \item \textbf{Unique Subsystem Contribution}: Unified Master Reference Architecture formalizing Constraint-First Architectural Synthesis (CFAS), invariant namespace (INV-01 to INV-15), and Theorem-Implementation Lattice.
    \item \textbf{Novelty Claim Formulation}: Portfolio-level meta-architecture reference model for privacy-first intelligent campus systems.
    \item \textbf{Theoretical Formalisms}: No formal mathematical theorem blocks.
\end{itemize}

\subsection*{Empirical Evidence & Telemetry}
\begin{itemize}[leftmargin=1.5em, itemsep=0.2em]
    \item \textbf{Key Numerical Telemetry}: Documented in manuscript.
    \item \textbf{Remaining Reviewer-6 Risk}: Survey/reference model format (bibliography completely rebuilt with external peer-reviewed literature).
\end{itemize}

\section{Paper 21: Formal Foundations of Spatiotemporal Compliance and Distributed System Integrity}

\subsection*{Paper Identity \& Strategic Metadata}
\begin{itemize}[leftmargin=1.5em, itemsep=0.2em]
    \item \textbf{Paper Identifier}: \texttt{P21} $\cdot$ \textbf{Methodological Type}: THEORETICAL (TOPOLOGICAL COMPLIANCE LOGIC)
    \item \textbf{Research Area}: Theoretical Computer Science / Formal Methods
    \item \textbf{Publication Status}: \textbf{DRAFT}
    \item \textbf{Target Venues}: Primary: \textit{ACM Transactions on Computational Logic} $\cdot$ Alternatives: \textit{Formal Methods in System Design, IEEE Trans. Automatic Control}
\end{itemize}

\subsection*{Scientific & Architectural Contribution Decomposition}
\begin{itemize}[leftmargin=1.5em, itemsep=0.2em]
    \item \textbf{Known Primitives Acknowledged}: Lebesgue measure theory, metric spaces, timed automata (Alur-Dill), Metric Temporal Logic.
    \item \textbf{Unique Subsystem Contribution}: Rigorous topological and measure-theoretic foundations proving Borel measurability, Nyquist completeness, and PSPACE-completeness of continuous compliance.
    \item \textbf{Novelty Claim Formulation}: 14 formal definitions and 8 foundational mathematical theorems with complete deductive proofs.
    \item \textbf{Theoretical Formalisms}: 22 formal mathematical blocks (theorem, definition)
\end{itemize}

\subsection*{Empirical Evidence & Telemetry}
\begin{itemize}[leftmargin=1.5em, itemsep=0.2em]
    \item \textbf{Key Numerical Telemetry}: Documented in manuscript.
    \item \textbf{Remaining Reviewer-6 Risk}: PSPACE complexity vs real-time execution (clarified as discrete polynomial-time PCVF approximation).
\end{itemize}

\section{Paper 22: Perception Integrity Foundations: Evidential Uncertainty, Disagreement Dynamics, and Blur Bounds in Edge Vision}

\subsection*{Paper Identity \& Strategic Metadata}
\begin{itemize}[leftmargin=1.5em, itemsep=0.2em]
    \item \textbf{Paper Identifier}: \texttt{P22} $\cdot$ \textbf{Methodological Type}: THEORETICAL \& EMPIRICAL (EVIDENTIAL PERCEPTION GATE)
    \item \textbf{Research Area}: Edge Computer Vision / Uncertainty Quantification
    \item \textbf{Publication Status}: \textbf{DRAFT}
    \item \textbf{Target Venues}: Primary: \textit{IEEE Transactions on Pattern Analysis and Machine Intelligence} $\cdot$ Alternatives: \textit{IEEE Trans. Neural Networks and Learning Systems, International Journal of Computer Vision}
\end{itemize}

\subsection*{Scientific & Architectural Contribution Decomposition}
\begin{itemize}[leftmargin=1.5em, itemsep=0.2em]
    \item \textbf{Known Primitives Acknowledged}: Evidential Deep Learning (Sensoy), Dirichlet distributions, MC-Dropout, Modified Laplacian blur.
    \item \textbf{Unique Subsystem Contribution}: Single-pass Layer-1 perception integrity gate with analytical Dirichlet variance upper bounds.
    \item \textbf{Novelty Claim Formulation}: Theorem 1 (Dirichlet Variance Upper Bound) and bounded composite risk $R_p \in [0,1]$.
    \item \textbf{Theoretical Formalisms}: 4 formal mathematical blocks (proposition, theorem, corollary)
\end{itemize}

\subsection*{Empirical Evidence & Telemetry}
\begin{itemize}[leftmargin=1.5em, itemsep=0.2em]
    \item \textbf{Key Numerical Telemetry}: AUROC=1.0000 / FPR95=0.0000 on 2,000-frame benchmark / ECE=0.0412 in 1.486 ms (Source: Physical ARM64 Dirichlet evidential inference over 5 corruption regimes)
    \item \textbf{Remaining Reviewer-6 Risk}: Synthetic corruption benchmark (addressed and defended in Section IV-C3/V-A).
\end{itemize}

\section{Paper 23: Adaptive Trustworthy Edge Systems: Dynamic Risk-Driven Cascades and Real-Time SLA Bounds}

\subsection*{Paper Identity \& Strategic Metadata}
\begin{itemize}[leftmargin=1.5em, itemsep=0.2em]
    \item \textbf{Paper Identifier}: \texttt{P23} $\cdot$ \textbf{Methodological Type}: THEORETICAL \& SYSTEMS (CASCADE SLA DUALITY)
    \item \textbf{Research Area}: Edge Computing / Real-Time Systems
    \item \textbf{Publication Status}: \textbf{DRAFT}
    \item \textbf{Target Venues}: Primary: \textit{IEEE Transactions on Mobile Computing} $\cdot$ Alternatives: \textit{IEEE Trans. Computers, ACM Trans. Embedded Computing Systems}
\end{itemize}

\subsection*{Scientific & Architectural Contribution Decomposition}
\begin{itemize}[leftmargin=1.5em, itemsep=0.2em]
    \item \textbf{Known Primitives Acknowledged}: BranchyNet, MSDNet, Fenchel-Rockafellar duality, M/G/1 queueing theory.
    \item \textbf{Unique Subsystem Contribution}: Dynamic risk-driven cascade architecture containing tail latencies (P99=4.556 ms) within a 5.0 ms SLA at 373.3 FPS capacity.
    \item \textbf{Novelty Claim Formulation}: Theorem 1 (Convex Continuum Duality) and Pollaczek-Khinchine delay bounds.
    \item \textbf{Theoretical Formalisms}: 2 formal mathematical blocks (proposition, theorem)
\end{itemize}

\subsection*{Empirical Evidence & Telemetry}
\begin{itemize}[leftmargin=1.5em, itemsep=0.2em]
    \item \textbf{Key Numerical Telemetry}: 373.3 FPS instantaneous capacity (2.679 ms) / P99=4.556 ms < 5.0 ms SLA / 8.1\% duty cycle (Source: Physical ARM64 M/G/1 queueing dispatch across 2,000 video inferences)
    \item \textbf{Remaining Reviewer-6 Risk}: Energy index proportional to FLOPs vs physical Joules (defended in Section III-C).
\end{itemize}

\section{Paper 24: Generalized Cross-Modal Recovery under Compromised Primary Sensing}

\subsection*{Paper Identity \& Strategic Metadata}
\begin{itemize}[leftmargin=1.5em, itemsep=0.2em]
    \item \textbf{Paper Identifier}: \texttt{P24} $\cdot$ \textbf{Methodological Type}: THEORETICAL \& EMPIRICAL (CROSS-MODAL JSD FUSION)
    \item \textbf{Research Area}: Multisensor Fusion / Information Theory
    \item \textbf{Publication Status}: \textbf{DRAFT}
    \item \textbf{Target Venues}: Primary: \textit{IEEE Transactions on Signal Processing} $\cdot$ Alternatives: \textit{IEEE Trans. Information Forensics and Security, IEEE Trans. Sensor Networks}
\end{itemize}

\subsection*{Scientific & Architectural Contribution Decomposition}
\begin{itemize}[leftmargin=1.5em, itemsep=0.2em]
    \item \textbf{Known Primitives Acknowledged}: Jensen-Shannon Divergence (JSD), Covariance Intersection, software Phase-Locked Loops (PLL).
    \item \textbf{Unique Subsystem Contribution}: Cross-modal consensus recovery with dynamic trust gradients transferring 95.0\% authority to intact secondary channels under 80\% optical degradation.
    \item \textbf{Novelty Claim Formulation}: Dynamic trust weight self- and cross-gradients under bounded JSD divergence and multi-rate ring buffer synchronization.
    \item \textbf{Theoretical Formalisms}: 4 formal mathematical blocks (proposition, theorem, definition, corollary)
\end{itemize}

\subsection*{Empirical Evidence & Telemetry}
\begin{itemize}[leftmargin=1.5em, itemsep=0.2em]
    \item \textbf{Key Numerical Telemetry}: Collapse of RGB trust (0.4000 to 0.0500) / Authority transfer to 95.0\% (0.4750 each) under 80\% noise (Source: Multi-rate ring buffer cross-modal evaluation over 2,000 frames)
    \item \textbf{Remaining Reviewer-6 Risk}: Secondary channels assumed clean in baseline trials (bounds analyzed in Section V-C3/VI).
\end{itemize}

\section{Paper 25: ScholarMaster Macro Integration Architecture and Downstream Error Propagation Analysis}

\subsection*{Paper Identity \& Strategic Metadata}
\begin{itemize}[leftmargin=1.5em, itemsep=0.2em]
    \item \textbf{Paper Identifier}: \texttt{P25} $\cdot$ \textbf{Methodological Type}: THEORETICAL \& SYSTEMS (MACRO ERROR PROPAGATION \& EAF)
    \item \textbf{Research Area}: Software Engineering / Cyber-Physical Safety
    \item \textbf{Publication Status}: \textbf{DRAFT}
    \item \textbf{Target Venues}: Primary: \textit{IEEE Transactions on Software Engineering} $\cdot$ Alternatives: \textit{IEEE Trans. Dependable and Secure Computing, ACM Trans. Cyber-Physical Systems}
\end{itemize}

\subsection*{Scientific & Architectural Contribution Decomposition}
\begin{itemize}[leftmargin=1.5em, itemsep=0.2em]
    \item \textbf{Known Primitives Acknowledged}: Data Cascades (Sambasivan), fault containment (Avizienis/Leveson), Voronoi diagrams, ArcFace angular margins.
    \item \textbf{Unique Subsystem Contribution}: Metric-space proof that nearest-neighbor biometric retrieval on hyperspheres exhibits essential step jump discontinuities (>=0.9589) across Voronoi facets, explaining macro error compounding and proving admitted-path EAF=0.0000 via Layer-1 quarantine.
    \item \textbf{Novelty Claim Formulation}: Theorem 1 (Voronoi Step Discontinuity), Proposition 1 (ArcFace Margin Lower Bound), and Proposition 2 (Piecewise Lipschitz Chain Rule).
    \item \textbf{Theoretical Formalisms}: 4 formal mathematical blocks (proposition, theorem, definition)
\end{itemize}

\subsection*{Empirical Evidence & Telemetry}
\begin{itemize}[leftmargin=1.5em, itemsep=0.2em]
    \item \textbf{Key Numerical Telemetry}: Admitted-path EAF=0.0000 via Layer-1 quarantine (78.4\% pass, 21.6\% quarantine) vs unprotected EAF=1.4220 (Source: 5-layer macro cascade evaluation across 2,000 evaluations under 5 corruption regimes)
    \item \textbf{Remaining Reviewer-6 Risk}: EAF=0.0000 achieved via fail-closed frame dropping (defended in Section IV-B/V-C).
\end{itemize}

\section{Single-Owner Contribution Lattice}
Every major scientific claim across the 25-paper portfolio is mapped to exactly one owning paper.

\section{Salami-Slicing and Methodological Overlap Analysis}
All closely related paper pairs have distinct research questions, methodologies, and evidence.

\section{Experimental Evidence Taxonomy and Matrix}
Empirical evidence is strictly categorized across Physical Measurements, Simulation Harnesses, Analytical Models, and User Studies.

\section{Hardware Strategy and Deployment Roadmap}
Maintains strict truthfulness regarding physical ARM64 deployments vs host simulation.

\section{Target Venue Strategy and Submission Positioning}
Venues are selected based on disciplinary alignment and transactions-level rigor.

\section{Publication-State Propagation Architecture}
Reusable multi-stage protocol for propagating newly published works into companion manuscripts without blind automatic insertion.

\section{Reusable Future-Paper Governance Template}
Standardized 12-point submission template for prospective ScholarMaster papers.

\section{Final Strategic Assessment and Freeze Recommendation}
Overall Portfolio Status: CONDITIONAL (Companion dependency alerts or unverified items pending review).

\newpage
\appendix
\section{Complete Portfolio Status Matrix}
All 25 manuscripts indexed with publication status and venue targets.

\end{document}
